\documentclass{amsart}
\usepackage{amsmath,amssymb,amsthm,mathtools}
\title{Homework 4}
\author{Antony Perez}
\begin{document}
\maketitle
\subsection*{7.1} Let $X$ and $Y$ be subspaces of a vector space $V$. Prove that $X \cap Y$ is a subspace of $V$. 
\begin{proof}
    Let $T = X \cap Y$. We know $T \neq \emptyset$, since both $X,Y$ contain the zero vector, due to their vector space properties. Let $x,y \in T$. Since $x,y \in X$, by the vector space properties, $x+y \in X$. Similarly, since $x,y \in Y$, $x+y \in Y$. Therefore, vector addition is in $T$. Let $\alpha \in \mathbb{F}$ be a scalar. Since $x \in X$, by vector space properties, $\alpha x \in X$. Similarly, $x \in Y \implies \alpha x \in Y$. Therefore, $\alpha x \in T$. Thus, $X \cap Y$ is a vector space.  
\end{proof}
\subsection*{7.2} Let $V$ be a vector space. For $X,Y \subset V$ the sum $X+Y$ is the collection of all vectors $v$ which can be represented as $v = x+y, x\in X, y \in Y$. Show that if $X$ and $Y$ are subspaces of $V$, then $X+Y$ is also a subspace. 
\begin{proof}
    Let $W = X+Y$. $W$ is non-empty, because both $X,Y$ contain the zero vector, due to their subspace properties ($0 + 0 = w \implies w = 0$). Let $w \in W$ such that for $x \in X, y \in Y$, $w = x + y$. Take two arbitrary vectors in $W$, named $w_1,w_2$ such that $w_1 = x_1 + y_1$ and $w_2 = x_2 + y_2$ for $x_1,x_2 \in X$ and $y_1, y_2 \in Y$. Adding both vectors: $w_1 + w_2 = x_1 + x_2 + y_1 + y_2$, where we see $x_1 + x_2 \in X$ and $y_1 + y_2 \in Y$ due to $X,Y$ being vector spaces. Thus, vector addition in $W$ is defined. Take an arbitrary scalar $\alpha \in \mathbb{F}$. Now let $\alpha w = \alpha x + \alpha y$. $\alpha x \in X$ and $\alpha y \in Y$ due to vector space properties. Thus, $\alpha w$ is defined, and $W$ has scalar multiplication. Therefore, $W$ is a vector space. 
\end{proof}
\subsection*{7.3} Let $X$ be a subspace of a vector space $V$, and let $v \in V, v \not \in X$. Prove that if $x \in X$, then $x+v \not \in X$.
\begin{proof}
    Suppose $x+v \in X$. Let $y \in X$ such that $x+v = y$. We may rewrite this as $v = y-x$. $y-x \in X$, due to vector addition and scalar multiplication (vector space properties). But $v \not \in X$. Contradiction. Therefore, $x+v \not \in X$
\end{proof}
\subsection*{7.4} Let $X$ and $Y$ be subspaces of a vector space $V$. Using the previous exercise, show that $X \cup Y$ is a subspace if and only if $X \subset Y$ or $Y \subset X$. 
\begin{proof}
    Suppose $X \not \subset Y$ and $Y \not \subset X$. Let $X \cup Y$ be a subspace, and let $x,y \in X \cup Y$ such that $x \in X$ but $x \not \in Y$, and $y \in Y$ but $y \not \in X$. Since $X \cup Y$ is a vector space, let $w \in X \cup Y$ be an arbitrary vector such that $w = x+y$. If $w \in X$, then $w - x = y$. $w-x \in X$ due to vector space properties, but $y \not \in X$. Contradiction. If $w \in Y$, then $w - y = x$. $w-y \in Y$ due to vector space properties, but $x \not \in Y$. Contradiction. Therefore, $X\cup Y$ is a vector space iff $X \subset Y$ or $Y \subset X$. Conversely, if $X \subset Y$, then $X\cup Y = Y$, and $Y$ is a vector space, so $X \cup Y$ is a vector space. if $Y \subset X$, then $X\cup Y = X$, and $X$ is a vector space, so $X \cup Y$ is a vector space. If $X \subset Y$ and $Y \subset X$, then $X = Y$, and $X \cup X = Y \cup Y$, and is automatically a vector space. Therefore, we have shown that $X\cup Y$ is a vector space iff $X \subset Y$ or $Y \subset X$. 
\end{proof}
\subsection*{2.1} Write the systems of equations below in matrix form and as vector equations:
\[
\text{a)} \quad
\left\{
\begin{aligned}
x_1 &+ 2x_2 &- x_3 &= -1 \\
2x_1 &+ 2x_2 &+ x_3 &= 1 \\
3x_1 &+ 5x_2 &- 2x_3 &= -1
\end{aligned}
\right.
\]

\[
\text{b)} \quad
\left\{
\begin{aligned}
x_1 &- 2x_2 &- x_3  &= 1 \\
2x_1 &- 3x_2 &+ x_3  &= 6 \\
3x_1 &- 5x_2 & & &= 7 \\
x_1 & & &+ 5x_3 &= 9
\end{aligned}
\right.
\]
Solve the systems. Write the answers in the vector form.
\begin{proof}
    For the first system: $R_2 \to R_2 - 2R_1 ;
R_3 \to R_3 - 3R_1 ;
R_2 \to -\tfrac{1}{2}R_2 ;
R_3 \to R_3 + R_2 ;
R_3 \to -2R_3 ;
R_1 \to R_1 + R_3 ;
R_2 \to R_2 + \tfrac{3}{2}R_3 ;
R_1 \to R_1 - 2R_2$ which gives $\begin{bmatrix}
    1 & 0 & 0 & 4 \\ 
    0 &1&0&-3 \\ 
    0&0&1&-1
\end{bmatrix}$. Its vector form is $4 \begin{bmatrix}
    1 \\ 2 \\ 3
\end{bmatrix} - 3\begin{bmatrix}
    2 \\ 2 \\ 5
\end{bmatrix} - \begin{bmatrix}
    -1 \\ 1 \\ -2
\end{bmatrix} = \begin{bmatrix}
    -1 \\ 1 \\ -1
\end{bmatrix}$ 

For the second system: $R_2 -= 2R_1; R_3 -= 3R_1; R_4 -= R_1; R_4 = 1/2 R_4; R_3 -= R_2; R_4 -= R_2; R_1 += 2R_2$, which gives the matrix $\begin{bmatrix}
    1&0&5&9 \\ 
    0&1&3&4 \\ 
    0&0&0&0 \\ 
    0&0&0&0 \\ 
\end{bmatrix}$. This shows $x_3$ is a free variable, so let $t = x_3$. Rewriting, we get $x_1 = 9-5t$, $x_2 = 4-3t, x_3 = t$. We may write the vector form as $\begin{bmatrix}
    9\\
    4\\
    0\\
\end{bmatrix} + t\begin{bmatrix}
    -5 \\ -3 \\ 1
\end{bmatrix} = b$
\end{proof}
\subsection*{3.1}
For what value of $b$ the system
\[
\begin{pmatrix}
1 & 2 & 2 \\
2 & 4 & 6 \\
1 & 2 & 3
\end{pmatrix}
\mathbf{x} =
\begin{pmatrix}
1 \\ 4 \\ b
\end{pmatrix}
\]
has a solution. Find the general solution of the system for this value of $b$.
\begin{proof}
    From RREF: $R_2 = 1/2R_2; R_2 -= R_1; R_3 -= R_1$, and since $x_3 = 1$, and $x_3 = b-1$, then $b=2$. Continuing on, $R_3 -= R_2; R_1 -2R_2$. $x_2$ is free, let $x_2 = t$. Thus, $x_1 = -1-2t, x_2 = t, x_3 = 1$. The general solution is $\begin{bmatrix}
        -1 \\ 0 \\ 1 
    \end{bmatrix} + t \begin{bmatrix}
        -2 \\ 1 \\ 0
    \end{bmatrix}$
\end{proof}
\subsection*{3.2}
Determine, if the vectors
\[
\begin{pmatrix} 1\\1\\0\\0 \end{pmatrix}, \quad
\begin{pmatrix} 1\\0\\1\\0 \end{pmatrix}, \quad
\begin{pmatrix} 0\\0\\1\\1 \end{pmatrix}, \quad
\begin{pmatrix} 0\\1\\0\\1 \end{pmatrix}
\]
are linearly independent or not.

Do these four vectors span $\mathbb{R}^4$? (In other words, is it a generating system?) What about $\mathbb{C}^4$?

\begin{proof}
    Augment the vectors with the zero vector such that $\begin{bmatrix}
        1 &1&0&0& 0\\ 
        1&0&0&1&0 \\ 
        0&1&1&0&0 \\ 
        0&0&1&1&0
    \end{bmatrix}$
    Doing RREF: $R_2 -= R_1; R_1 += R_2; R_3 += R_2; R_2 = -R_2; R_4 -=R_3$. Since there are free variables, this is not linearly independent. So, it cannot generate $\mathbb{R}^4$. Same with $\mathbb{C}^4$, for the same reasons. 
\end{proof}
\subsection*{3.3}
Determine, which of the following systems of vectors are bases in $\mathbb{R}^3$:
\begin{enumerate}
\item[a)] $(1,2,-1)^T$, $(1,0,2)^T$, $(2,1,1)^T$;
\item[b)] $(-1,3,2)^T$, $(-3,1,3)^T$, $(2,10,2)^T$;
\item[c)] $(67,13,-47)^T$, $(\pi,-7.84,0)^T$, $(3,0,0)^T$.
\end{enumerate}
Which of the systems are bases in $\mathbb{C}^3$?

\begin{proof}
    From the first system, we use RREF to show linear independence due to pivots in every column and row: $R_2 += 3R_3; R_3 += R_1; R_2 -= R_3; R_3 = 1/2 R_3; R_1 -= R_2; R_3 -= R_2; R_1 -= 2R_2$. Thus, since it is a $3\times 3$ matrix and it is linearly independent, the vectors span $\mathbb{R}^3$ and $\mathbb{C}^3$, since $\mathbb{R} \subset \mathbb{C}$.
    
    For the second system, $R_2 += 3R_1; R_3 += 2R_1; R_2 = 1/8 R_2; R_3 = 1/3 R_3$. Notice both $R_3$ and $R_2$ are identical. Thus, we are guaranteed a free variable. So, system 2 cannot be a basis for $\mathbb{R}^3$ nor $\mathbb{C}^3$.

    Lastly, the weird system. Notice if we $R_1 \leftrightarrow R_3$, we get an lower triangle, with non-zero pivots. Thus, this system is linearly independent and spans both $\mathbb{C}^3$ and $\mathbb{R}^3$
\end{proof}
\subsection*{3.4}
Do the polynomials $x^3+2x$, $x^2+x+1$, $x^3+5$ generate (span) $\mathbb{P}_3$? Justify your answer.
\begin{proof}
    No. This is because there are only 3 columns, where $\mathbb{P}_4$ requires 4 columns, because of the need for an $1,x,x^2,$ and $x^3$. 
\end{proof}

\subsection*{3.5} Can 5 vectors in $\mathbb{F}^4$ be linearly independent? Justify your answer.
\begin{proof}
    No. Since there are 5 vectors, you will have a $4 \times 5$ matrix. Since there are more columns than rows, you will get a free variable. 
\end{proof}

\subsection*{3.6}
Prove or disprove: If the columns of a square ($n\times n$) matrix $A$ are linearly independent, so are the columns of $A^2 = AA$.
\begin{proof}
    Suppose $A^2x=0$. Let $Ax = b$, where $x$ is some vector. Multiplying both sides, we get $A^2x = Ab=0$. Since $A$ is linearly independent, $Ax = 0 \implies b=0$. Since the columns are linearly independent, $Ax = 0 \implies x = 0$. Therefore, $A^2x = 0$, showing the columns are linearly independent. 
\end{proof}

\subsection*{3.7}  Prove or disprove: If the columns of a square ($n\times n$) matrix $A$ are linearly independent, so are the rows of $A^3 = AAA$.
\begin{proof}
    Using 3.6, we can apply the same technique and get $A^3x = 0$, and applying linearity, we see that $A^3x = 0 \implies x = 0$. Since $A^3$ is square, columns being independent $\Longleftrightarrow$ independent rows. So the rows of $A^3$ are independent as well. 
\end{proof}

\end{document}